% 康奈尔听课笔记模板（配色取自王道课件）
\documentclass[10.5pt,a4paper,oneside]{article}
\usepackage[a4paper,top=2.3cm,bottom=2.2cm,left=5.7cm,right=1.5cm,
  marginparwidth=4.25cm,marginparsep=0.55cm,headheight=14pt,footskip=1.1cm]{geometry}
\usepackage{xeCJK}
\usepackage{xeCJKfntef}
\xeCJKDeclareCharClass{FullLeft}{"201C,"2018}
\xeCJKDeclareCharClass{FullRight}{"201D,"2019}
\xeCJKDeclareCharClass{CJK}{"2460->"24FF,"2605,"2606}
\setmainfont{Noto Sans CJK SC}[BoldFont={Noto Sans CJK SC Bold},ItalicFont={Noto Sans CJK SC},ItalicFeatures={FakeSlant=0.18}]
\setsansfont{Noto Sans CJK SC}[BoldFont={Noto Sans CJK SC Bold}]
\setmonofont{DejaVu Sans Mono}[Scale=0.9]
\setCJKmainfont{Noto Sans CJK SC}[BoldFont={Noto Sans CJK SC Bold},ItalicFont={Noto Sans CJK SC},ItalicFeatures={FakeSlant=0.18}]
\setCJKsansfont{Noto Sans CJK SC}[BoldFont={Noto Sans CJK SC Bold}]
\setCJKmonofont{Noto Sans Mono CJK SC}[Scale=0.9]
\linespread{1.32}
\usepackage{amsmath,amssymb}
\usepackage[dvipsnames,table]{xcolor}
\usepackage{graphicx}
\usepackage{adjustbox}
\usepackage{tikz}
\usetikzlibrary{calc}
\usepackage[most]{tcolorbox}
\usepackage{eso-pic}
\usepackage{fancyhdr}
\usepackage{titlesec}
\usepackage{needspace}
\usepackage{enumitem}
\usepackage{xltabular,booktabs,array,colortbl}
\usepackage{marginnote}
\reversemarginpar
\renewcommand*{\raggedleftmarginnote}{}
\renewcommand*{\raggedrightmarginnote}{}
\usepackage{caption}
\usepackage{titletoc}
\usepackage{fancyvrb}
\usepackage{framed}
\usepackage[section]{placeins}
\usepackage{hyperref}
\hypersetup{colorlinks=true,linkcolor=wdblue,urlcolor=wdblue,pdftitle={3.2.2
队列的顺序实现},pdfauthor={听课笔记}}

% ---------- 课件配色 ----------
\definecolor{wdblue}{HTML}{4472C4}
\definecolor{wdlightblue}{HTML}{5B9BD5}
\definecolor{wdtitle}{HTML}{4A86C8}
\definecolor{wddark}{HTML}{1F3864}
\definecolor{wdorange}{HTML}{ED7D31}
\definecolor{wdgray}{HTML}{A5A5A5}
\definecolor{wdgold}{HTML}{FFC000}
\definecolor{wdgreen}{HTML}{70AD47}
\definecolor{wddeepgreen}{HTML}{3F7F2A}
\definecolor{wdred}{HTML}{C00000}
\definecolor{wdpurple}{HTML}{7030A0}
\definecolor{cuebg}{HTML}{F2F6FC}
\definecolor{codebg}{HTML}{F7F9FC}
\color[HTML]{262626}

% ---------- 行内样式 ----------
\newcommand{\hlred}[1]{\textcolor{wdred}{\bfseries #1}}
\newcommand{\hlblue}[1]{\textcolor{wdblue}{\bfseries #1}}
\newcommand{\hlgreen}[1]{\textcolor{wddeepgreen}{\bfseries #1}}
\newcommand{\hlorange}[1]{\textcolor{wdorange}{\bfseries #1}}
\newcommand{\hlpurple}[1]{\textcolor{wdpurple}{\bfseries #1}}
\newcommand{\uln}[1]{\CJKunderline[format=\color{wdred},thickness=0.9pt,depth=0.28em]{#1}}
\newcommand{\wavy}[1]{\CJKunderwave[format=\color{wdorange}]{#1}}
\newcommand{\hlmark}[1]{\CJKunderline[format=\color{wdgold!55},thickness=0.55em,depth=-0.05em]{#1}}
\newcommand{\tagbox}[2]{\tcbox[on line,colback=#1,colframe=#1,coltext=white,boxsep=0.5pt,left=3pt,right=3pt,top=0.5pt,bottom=0.5pt,arc=2pt,fontupper=\small\bfseries]{#2}}

% ---------- 康奈尔版式背景：左侧线索栏 ----------
\newlength{\cuestripw}\setlength{\cuestripw}{5.15cm}
\AddToShipoutPictureBG{%
  \begin{tikzpicture}[remember picture,overlay]
    \fill[cuebg] (current page.north west) rectangle ($(current page.south west)+(\cuestripw,0)$);
    \draw[wdblue,line width=1.2pt] ($(current page.north west)+(\cuestripw,-1.55cm)$) -- ($(current page.south west)+(\cuestripw,1.55cm)$);
    \node[anchor=west,text=wdblue,font=\small\bfseries] at ($(current page.north west)+(0.75cm,-1.25cm)$) {线索 · 问题};
    \node[anchor=west,text=wdblue,font=\small\bfseries] at ($(current page.north west)+(\cuestripw+0.55cm,-1.25cm)$) {笔记 · 老师主线 + 补充讲解};
  \end{tikzpicture}}

% ---------- 页眉页脚 ----------
\pagestyle{fancy}\fancyhf{}
\renewcommand{\headrulewidth}{0pt}
\fancyhead[R]{\small\color{wdgray}第3章 栈、队列和数组}
\fancyfoot[C]{\small\color{wdgray}\thepage}
\fancyfoot[R]{\small\color{wdorange}3.2.2 队列的顺序实现 · 听课笔记}
\fancyfoot[L]{\small\hyperlink{toc}{\textcolor{wdblue}{\ensuremath{\uparrow}\,回到目录}}}
\fancypagestyle{plain}{\fancyhf{}\fancyfoot[C]{\small\color{wdgray}\thepage}\fancyfoot[R]{\small\color{wdorange}3.2.2
队列的顺序实现 · 听课笔记}}

% ---------- 标题：仿课件的蓝色标题 + 六色条 ----------
\newcommand{\sixbar}{\par\nointerlineskip\vspace{3pt}\centering
  \begin{tikzpicture}\foreach \c [count=\i] in {wdlightblue,wdorange,wdgray,wdgold,wdblue,wdgreen}
     {\fill[\c] ({(\i-1)*1.05},0) rectangle ({(\i-1)*1.05+0.92},0.11);}\end{tikzpicture}\par}
\titleformat{\section}[block]{\centering\Large\bfseries\color{wdtitle}}{}{0pt}{}[\sixbar]
\titlespacing*{\section}{0pt}{16pt}{10pt}
\pretocmd{\section}{\Needspace{8\baselineskip}}{}{}
\titleformat{\subsection}{\Needspace{7\baselineskip}\large\bfseries\color{wddark}}{\textcolor{wdblue}{\rule[-0.15em]{4pt}{1.05em}}\hspace{6pt}}{0pt}{}
\titlespacing*{\subsection}{0pt}{12pt}{5pt}
\titleformat{\subsubsection}{\Needspace{5\baselineskip}\normalsize\bfseries\color{wdblue}}{\textcolor{wdorange}{\large\ensuremath{\blacktriangleright}}\hspace{4pt}}{0pt}{}
\titlespacing*{\subsubsection}{0pt}{8pt}{3pt}
\setcounter{secnumdepth}{0}
\setlength{\parindent}{0pt}\setlength{\parskip}{5pt}
\setlist{nosep,leftmargin=1.5em,itemsep=2pt}
\setlist[itemize,1]{label=\textcolor{wdblue}{\small\ensuremath{\bullet}}}
\setlist[itemize,2]{label=\textcolor{wdorange}{\scriptsize\ensuremath{\blacktriangleright}}}
\setlist[enumerate,1]{label=\textcolor{wdblue}{\bfseries\arabic*.}}
\providecommand{\tightlist}{\setlength{\itemsep}{2pt}\setlength{\parskip}{0pt}}

% ---------- 目录 ----------
\setcounter{tocdepth}{1}
\titlecontents{section}[0pt]{\addvspace{5pt}\bfseries\color{wddark}}{}{}{\titlerule*[0.7pc]{\color{wdgray!60}.}\contentspage}
\titlecontents{subsection}[1.6em]{\addvspace{1pt}\small\color{black!75}}{}{}{\titlerule*[0.7pc]{\color{wdgray!40}.}\contentspage}
\newcommand{\tocback}{\par\vspace{-3pt}{\hfill\footnotesize\hyperlink{toc}{\textcolor{wdgray}{\ensuremath{\uparrow}\,回到目录}}}\par\vspace{-4pt}}

% ---------- 线索栏 ----------
\newcommand{\cueblock}[2]{\Needspace{4\baselineskip}\leavevmode#1{#2}\par\vspace{-\baselineskip}\vspace{-\parskip}}
% 线索直接贴在段首行的左侧（不用 marginnote，避免分页时偶尔落进正文栏）
\newcommand{\cuebox}[1]{\leavevmode\llap{\smash{\begin{minipage}[t]{\marginparwidth}\raggedright\small\linespread{1.18}\selectfont #1\end{minipage}}\hspace{\marginparsep}}}
\newcommand{\cuenote}[1]{\cuebox{{\color{wdblue}\bfseries Q\,}\color{wddark}#1}}
\newcommand{\cuekey}[1]{\cuebox{\color{wdred}\bfseries #1}}

% ---------- 讲解框 ----------
\tcbset{wdbox/.style={enhanced,breakable,arc=2pt,boxrule=0pt,leftrule=3pt,left=7pt,right=7pt,top=5pt,bottom=5pt,
   before skip=7pt,after skip=7pt,fonttitle=\bfseries\small,attach title to upper={\quad},coltitle=#1,
   colframe=#1,colback=#1!6}}
\newtcolorbox{defbox}[1][]{enhanced,breakable,arc=3pt,boxrule=0.8pt,colframe=wdblue,colback=wdblue!4,
   left=8pt,right=8pt,top=6pt,bottom=6pt,fonttitle=\bfseries\small,coltitle=white,colbacktitle=wdblue,
   title={定义#1},before skip=8pt,after skip=8pt}
\newtcolorbox{thmbox}[1]{enhanced,breakable,arc=3pt,boxrule=0.8pt,colframe=wdblue,colback=wdblue!4,
   left=8pt,right=8pt,top=6pt,bottom=6pt,fonttitle=\bfseries\small,coltitle=white,colbacktitle=wdblue,
   title={#1},before skip=8pt,after skip=8pt}
\newtcolorbox{bridgebox}{wdbox=wdblue,title={承上启下}}
\newtcolorbox{plainbox}{wdbox=wddeepgreen,colframe=wdgreen,colback=wdgreen!8,title={大白话}}
\newtcolorbox{deepbox}{wdbox=wdpurple,title={背后的思想}}
\newtcolorbox{exambox}{wdbox=wdred,title={考点}}
\newtcolorbox{pitfallbox}{wdbox=wdorange,title={易错提醒}}
\newtcolorbox{bookbox}{wdbox=wddark,colback=wdgray!10,colframe=wdgray,title={王道书补充}}
\newtcolorbox{teacherbox}{wdbox=wdorange,colback=wdgold!8,colframe=wdgold,coltitle=wdorange,title={老师原话}}
\newtcolorbox{fixbox}{wdbox=wdred,colback=wdred!4,title={勘误 / 说明}}
% 习题笔记用
\newtcolorbox{questionbox}[1][]{enhanced,breakable,arc=3pt,boxrule=0.8pt,colframe=wddark,colback=white,
   left=8pt,right=8pt,top=6pt,bottom=6pt,fonttitle=\bfseries\small,coltitle=white,colbacktitle=wddark,
   title={#1},before skip=10pt,after skip=6pt}
\newtcolorbox{videobox}{wdbox=wdorange,colback=wdgold!8,colframe=wdgold,coltitle=wdorange,title={视频讲解}}
\newtcolorbox{jiexibox}{wdbox=wddark,colback=wdgray!10,colframe=wdgray,title={书上解析}}
\newtcolorbox{addbox}{wdbox=wdpurple,title={我的补充}}
\newtcolorbox{recallbox}{wdbox=wdblue,title={回顾正课}}
\newtcolorbox{watchbox}{enhanced,breakable,arc=3pt,boxrule=1pt,colframe=wdorange,colback=wdgold!12,
   fonttitle=\bfseries\small,coltitle=white,colbacktitle=wdorange,title={建议去看视频},left=8pt,right=8pt,before skip=8pt,after skip=8pt}
\newtcolorbox{summarybox}{enhanced,arc=3pt,
   boxrule=1pt,colframe=wdblue,colback=white,colbacktitle=wdblue,coltitle=white,fonttitle=\bfseries,
   title={总结区 · 用自己的话复述这一节},left=10pt,right=10pt,top=8pt,bottom=8pt,before skip=14pt}

% ---------- 代码 ----------
\usepackage{color}
\usepackage{fancyvrb}
\newcommand{\VerbBar}{|}
\newcommand{\VERB}{\Verb[commandchars=\\\{\}]}
\DefineVerbatimEnvironment{Highlighting}{Verbatim}{commandchars=\\\{\}}
% Add ',fontsize=\small' for more characters per line
\newenvironment{Shaded}{}{}
\newcommand{\AlertTok}[1]{\textcolor[rgb]{1.00,0.00,0.00}{#1}}
\newcommand{\AnnotationTok}[1]{\textcolor[rgb]{0.38,0.63,0.69}{#1}}
\newcommand{\AttributeTok}[1]{\textcolor[rgb]{0.49,0.56,0.16}{#1}}
\newcommand{\BaseNTok}[1]{\textcolor[rgb]{0.77,0.35,0.07}{#1}}
\newcommand{\BuiltInTok}[1]{\textcolor[rgb]{0.00,0.50,0.00}{#1}}
\newcommand{\CharTok}[1]{\textcolor[rgb]{0.63,0.19,0.43}{#1}}
\newcommand{\CommentTok}[1]{\textcolor[rgb]{0.18,0.55,0.23}{#1}}
\newcommand{\CommentVarTok}[1]{\textcolor[rgb]{0.18,0.55,0.23}{#1}}
\newcommand{\ConstantTok}[1]{\textcolor[rgb]{0.77,0.35,0.07}{#1}}
\newcommand{\ControlFlowTok}[1]{\textcolor[rgb]{0.12,0.31,0.71}{#1}}
\newcommand{\DataTypeTok}[1]{\textcolor[rgb]{0.12,0.31,0.71}{#1}}
\newcommand{\DecValTok}[1]{\textcolor[rgb]{0.77,0.35,0.07}{#1}}
\newcommand{\DocumentationTok}[1]{\textcolor[rgb]{0.18,0.55,0.23}{#1}}
\newcommand{\ErrorTok}[1]{\textcolor[rgb]{1.00,0.00,0.00}{#1}}
\newcommand{\ExtensionTok}[1]{#1}
\newcommand{\FloatTok}[1]{\textcolor[rgb]{0.77,0.35,0.07}{#1}}
\newcommand{\FunctionTok}[1]{\textcolor[rgb]{0.44,0.19,0.63}{#1}}
\newcommand{\ImportTok}[1]{\textcolor[rgb]{0.44,0.19,0.63}{#1}}
\newcommand{\InformationTok}[1]{\textcolor[rgb]{0.38,0.63,0.69}{#1}}
\newcommand{\KeywordTok}[1]{\textcolor[rgb]{0.12,0.31,0.71}{#1}}
\newcommand{\NormalTok}[1]{#1}
\newcommand{\OperatorTok}[1]{\textcolor[rgb]{0.20,0.20,0.20}{#1}}
\newcommand{\OtherTok}[1]{\textcolor[rgb]{0.00,0.44,0.13}{#1}}
\newcommand{\PreprocessorTok}[1]{\textcolor[rgb]{0.44,0.19,0.63}{#1}}
\newcommand{\RegionMarkerTok}[1]{#1}
\newcommand{\SpecialCharTok}[1]{\textcolor[rgb]{0.63,0.19,0.43}{#1}}
\newcommand{\SpecialStringTok}[1]{\textcolor[rgb]{0.63,0.19,0.43}{#1}}
\newcommand{\StringTok}[1]{\textcolor[rgb]{0.63,0.19,0.43}{#1}}
\newcommand{\VariableTok}[1]{\textcolor[rgb]{0.10,0.09,0.49}{#1}}
\newcommand{\VerbatimStringTok}[1]{\textcolor[rgb]{0.63,0.19,0.43}{#1}}
\newcommand{\WarningTok}[1]{\textcolor[rgb]{0.38,0.63,0.69}{#1}}
\makeatletter
\@ifundefined{Shaded}{}{\renewenvironment{Shaded}{\def\FrameCommand{\fboxsep=5pt\fboxrule=0.5pt\fcolorbox{wdblue!35}{codebg}}%
   \MakeFramed{\advance\hsize-\width\FrameRestore}\small}{\endMakeFramed}}
\makeatother
\fvset{fontsize=\small}

% ---------- 表格 ----------
\renewcommand{\arraystretch}{1.42}
\newcolumntype{Y}{>{\raggedright\arraybackslash}X}

% ---------- 图片 ----------
% 课件截图做成浮动体：LaTeX 会自动用正文填满页面，避免大块空白
\setcounter{topnumber}{3}\setcounter{bottomnumber}{2}\setcounter{totalnumber}{4}
\renewcommand{\topfraction}{0.85}\renewcommand{\bottomfraction}{0.6}\renewcommand{\textfraction}{0.12}\renewcommand{\floatpagefraction}{0.75}
\setlength{\textfloatsep}{10pt plus 3pt minus 2pt}\setlength{\floatsep}{8pt plus 2pt}
\newcommand{\slidefig}[3]{\begin{figure}[!ht]\centering\begin{tcolorbox}[enhanced,width=#2\linewidth,arc=2pt,boxrule=0.6pt,colframe=wdgray!70,
   colback=white,left=2pt,right=2pt,top=2pt,bottom=2pt,drop fuzzy shadow=black!25]\includegraphics[width=\linewidth]{#1}\end{tcolorbox}
   \vspace{-3pt}{\small\color{wdgray}#3}\end{figure}}
\newcommand{\plainfig}[3]{\par\noindent\begin{minipage}{\linewidth}\centering\includegraphics[width=#2\linewidth]{#1}\par\vspace{2pt}{\small\color{wdgray}#3}\end{minipage}\par\vspace{4pt}}

\begin{document}
% ---------- 封面条（仿课件封面） ----------
\thispagestyle{plain}
\begin{tcolorbox}[enhanced,enlarge left by=-4.8cm,width=\linewidth+4.8cm,arc=0pt,boxrule=0pt,
   colback=wddark,interior style={left color=wddark,right color=wdblue},left=18pt,right=18pt,top=14pt,bottom=12pt,
   overlay={\draw[white,line width=2.5pt] ([xshift=10pt,yshift=-10pt]frame.north west) -- ++(0,-0.9cm);
            \draw[white,line width=2.5pt] ([xshift=10pt,yshift=-10pt]frame.north west) -- ++(1.6cm,0);
            \draw[white,line width=2.5pt] ([xshift=-10pt,yshift=10pt]frame.south east) -- ++(0,0.9cm);
            \draw[white,line width=2.5pt] ([xshift=-10pt,yshift=10pt]frame.south east) -- ++(-1.6cm,0);}]
  {\color{white!85!wdlightblue}\small 数据结构 · 第3章
栈、队列和数组 · 康奈尔听课笔记}\par\vspace{6pt}
  {\color{white}\fontsize{22}{28}\selectfont\bfseries 3.2.2
队列的顺序实现}\par\vspace{6pt}
  {\color{white!80!wdlightblue}\small 课件 3.2.2 · 正课字幕 P24 · 王道书
§3.2.2（第 76--79 页）}
\end{tcolorbox}
\vspace{4pt}
\begin{tcolorbox}[enhanced,arc=3pt,boxrule=0.8pt,colframe=wdblue,colback=wdblue!3,left=10pt,right=10pt,top=6pt,bottom=6pt,
   colbacktitle=wdblue,coltitle=white,fonttitle=\bfseries\small,title={目录 · 点击标题跳转}]
\hypertarget{toc}{}\makeatletter\@starttoc{toc}\makeatother
\end{tcolorbox}
\hypertarget{ux8fd9ux4e00ux8282ux5728ux8bb2ux4ec0ux4e48}{%
\section{这一节在讲什么}\label{ux8fd9ux4e00ux8282ux5728ux8bb2ux4ec0ux4e48}}

\tocback

\Needspace{4\baselineskip}

\cuenote{用顺序存储实现队列，最难的是哪几件事？}用\textbf{静态数组}存放队列元素，再设两个变量标记队头和队尾。确定存储结构后，讨论这些基本操作怎么实现：创（初始化）、增（入队）、删（出队）、查（获取队头元素），以及进行增、删、查之前必要的\textbf{判空、判满}。

这一节\hlred{很容易在选择题中考察}，难点集中在三件事：怎么让空间“循环”起来、怎么区分队空和队满、怎么算队列长度。

\hypertarget{ux4e00ux987aux5e8fux961fux5217ux7684ux5b9aux4e49}{%
\section{一、顺序队列的定义}\label{ux4e00ux987aux5e8fux961fux5217ux7684ux5b9aux4e49}}

\tocback

\cueblock{\cuenote}{\relax{\texttt{front}、\texttt{rear} 各指向哪里？\texttt{SqQueue\ Q;}
分配了多大的空间？}}

\begin{Shaded}
\begin{Highlighting}[]
\PreprocessorTok{\#define MaxSize }\DecValTok{10}\PreprocessorTok{                  }\CommentTok{// 定义队列中元素的最大个数}
\KeywordTok{typedef} \KeywordTok{struct} \OperatorTok{\{}
\NormalTok{    ElemType data}\OperatorTok{[}\NormalTok{MaxSize}\OperatorTok{];}         \CommentTok{// 用静态数组存放队列元素}
    \DataTypeTok{int}\NormalTok{ front}\OperatorTok{,}\NormalTok{ rear}\OperatorTok{;}                \CommentTok{// 队头指针和队尾指针}
\OperatorTok{\}}\NormalTok{ SqQueue}\OperatorTok{;}                          \CommentTok{// Sq：sequence，顺序}

\DataTypeTok{void}\NormalTok{ testQueue}\OperatorTok{()} \OperatorTok{\{}
\NormalTok{    SqQueue Q}\OperatorTok{;}                      \CommentTok{// 声明一个队列（顺序存储）}
    \CommentTok{// ……后续操作}
\OperatorTok{\}}
\end{Highlighting}
\end{Shaded}

\begin{itemize}
\tightlist
\item
  执行 \texttt{SqQueue\ Q;} 后，系统分配一整片连续空间，大小
  \texttt{MaxSize*sizeof(ElemType)}，能存 10 个数据元素；
\item
  规定：\uln{\texttt{front} 指向队头元素；\texttt{rear}
  指向队尾元素的\textbf{后一个}位置}，也就是下一个应该插入数据元素的位置。课件例子里
  a、b、c、d、e 在队中（下面是队头，上面是队尾），\texttt{front} 是
  0，\texttt{rear} 是 5。
\end{itemize}

\slidefig{assets/3.2.2-p03.jpg}{0.6}{课件：队列的顺序实现（front 指向队头元素，rear
指向队尾元素的后一个位置）}

\hypertarget{ux4e8cux521dux59cbux5316ux4e0eux5224ux7a7a}{%
\section{二、初始化与判空}\label{ux4e8cux521dux59cbux5316ux4e0eux5224ux7a7a}}

\tocback

\Needspace{4\baselineskip}

\cuenote{初始化时 \texttt{front}、\texttt{rear} 为什么都设成
0？判空的条件是什么？}\texttt{rear}
指向接下来应该插入元素的位置，所以初始化时让队头、队尾指针\uln{都指向 0}。于是可以用两个指针是否相等来判断队列是否为空。

\begin{Shaded}
\begin{Highlighting}[]
\CommentTok{// 初始化队列}
\DataTypeTok{void}\NormalTok{ InitQueue}\OperatorTok{(}\NormalTok{SqQueue }\OperatorTok{\&}\NormalTok{Q}\OperatorTok{)} \OperatorTok{\{}
\NormalTok{    Q}\OperatorTok{.}\NormalTok{rear }\OperatorTok{=}\NormalTok{ Q}\OperatorTok{.}\NormalTok{front }\OperatorTok{=} \DecValTok{0}\OperatorTok{;}           \CommentTok{// 初始时 队头、队尾指针指向 0}
\OperatorTok{\}}

\CommentTok{// 判断队列是否为空}
\DataTypeTok{bool}\NormalTok{ QueueEmpty}\OperatorTok{(}\NormalTok{SqQueue Q}\OperatorTok{)} \OperatorTok{\{}
    \ControlFlowTok{if} \OperatorTok{(}\NormalTok{Q}\OperatorTok{.}\NormalTok{rear }\OperatorTok{==}\NormalTok{ Q}\OperatorTok{.}\NormalTok{front}\OperatorTok{)}          \CommentTok{// 队空条件}
        \ControlFlowTok{return} \KeywordTok{true}\OperatorTok{;}
    \ControlFlowTok{else}
        \ControlFlowTok{return} \KeywordTok{false}\OperatorTok{;}
\OperatorTok{\}}
\end{Highlighting}
\end{Shaded}

\slidefig{assets/3.2.2-p04.jpg}{0.6}{课件：初始化操作与判空}

\hypertarget{ux4e09ux5165ux961fux4e0eux5faaux73afux961fux5217}{%
\section{三、入队与循环队列}\label{ux4e09ux5165ux961fux4e0eux5faaux73afux961fux5217}}

\tocback

\Needspace{4\baselineskip}

\cuenote{\texttt{rear\ ==\ MaxSize}
能不能作为队满条件？\texttt{(Q.rear+1)\%MaxSize} 的作用是什么？}只能从\textbf{队尾}入队。静态数组容量有限，插入前要先判断队列是否已满（怎么判断，第五节再讨论）：

\begin{Shaded}
\begin{Highlighting}[]
\DataTypeTok{bool}\NormalTok{ EnQueue}\OperatorTok{(}\NormalTok{SqQueue }\OperatorTok{\&}\NormalTok{Q}\OperatorTok{,}\NormalTok{ ElemType x}\OperatorTok{)} \OperatorTok{\{}
    \ControlFlowTok{if} \OperatorTok{(}\NormalTok{队列已满}\OperatorTok{)}
        \ControlFlowTok{return} \KeywordTok{false}\OperatorTok{;}               \CommentTok{// 队满则报错}
\NormalTok{    Q}\OperatorTok{.}\NormalTok{data}\OperatorTok{[}\NormalTok{Q}\OperatorTok{.}\NormalTok{rear}\OperatorTok{]} \OperatorTok{=}\NormalTok{ x}\OperatorTok{;}             \CommentTok{// 将 x 插入队尾}
\NormalTok{    Q}\OperatorTok{.}\NormalTok{rear }\OperatorTok{=}\NormalTok{ Q}\OperatorTok{.}\NormalTok{rear }\OperatorTok{+} \DecValTok{1}\OperatorTok{;}            \CommentTok{// 队尾指针后移}
    \ControlFlowTok{return} \KeywordTok{true}\OperatorTok{;}
\OperatorTok{\}}
\end{Highlighting}
\end{Shaded}

照这个逻辑一直入队，存满后 \texttt{rear} 等于 10。那能不能认为
\texttt{rear\ ==\ MaxSize}
时队列就满了？\hlred{不能}：如果之后有几个队头元素出队，\texttt{front}
会往后移，数组前面这些位置就空出来了，还可以继续插入。所以
\texttt{rear\ ==\ MaxSize} 时队列并不一定满。

\slidefig{assets/3.2.2-p05.jpg}{0.6}{课件：入队操作；“队列已满的条件：rear==MaxSize ???”}

\begin{bookbox}

王道书第 77 页给这个现象起了名字：队列中仅剩一个元素，但 \texttt{rear}
已达
\texttt{MaxSize}，若此时继续入队，则会出现“\textbf{上溢}”；但这种溢出并非真正溢出，实际上数组中仍有空闲单元，因此称此现象为\hlred{\textbf{假溢出}}。

\end{bookbox}

\slidefig{assets/3.2-book-fig3.6.png}{0.6}{王道书图 3.6：队列的操作 (a) 空队 (b) 5 个元素入队 (c) 出队 1 次 (d)
出队 3 次}

\hypertarget{ux7528ux53d6ux6a21ux8ba9-rear-ux56deux5230ux5f00ux5934}{%
\subsection{用取模让 rear
回到开头}\label{ux7528ux53d6ux6a21ux8ba9-rear-ux56deux5230ux5f00ux5934}}

把第二句改成取余运算：

\begin{Shaded}
\begin{Highlighting}[]
\NormalTok{Q}\OperatorTok{.}\NormalTok{data}\OperatorTok{[}\NormalTok{Q}\OperatorTok{.}\NormalTok{rear}\OperatorTok{]} \OperatorTok{=}\NormalTok{ x}\OperatorTok{;}                 \CommentTok{// 新元素插入队尾}
\NormalTok{Q}\OperatorTok{.}\NormalTok{rear }\OperatorTok{=} \OperatorTok{(}\NormalTok{Q}\OperatorTok{.}\NormalTok{rear }\OperatorTok{+} \DecValTok{1}\OperatorTok{)} \OperatorTok{\%}\NormalTok{ MaxSize}\OperatorTok{;}    \CommentTok{// 队尾指针加 1 取模}
\end{Highlighting}
\end{Shaded}

设 \texttt{rear} 是 9（数组最后一个位置），新元素放进
\texttt{data{[}9{]}}
后，\texttt{rear\ =\ (9+1)\%10\ =\ 0}，又指回了数组开头。

\begin{addbox}

\textbf{跨考 Tips：取模运算即取余运算}。\texttt{a\ \%\ b} 等于 a 除以 b
的余数，在《数论》中通常写作
\texttt{a\ MOD\ b}。例：\texttt{23\ \%\ 7}，23÷7 = 3 余 2，所以结果是
2。任何整数除以 7，余数只可能是 0～6，不可能大于等于
7。所以模运算把\uln{无限的整数域映射到了有限的整数集合}
\(\{0, 1, 2, \dots, b-1\}\) 上。

\end{addbox}

这样 \texttt{rear}
从下到上、再回到最下面、再往上，循环往复。相当于把线状的存储空间，\hlred{在逻辑上变成了“环状”}。用这种方式实现的队列叫\textbf{循环队列}，这个术语可能在选择题中出现。

\slidefig{assets/3.2.2-p06.jpg}{0.6}{课件：用模运算将存储空间在逻辑上变成“环状”}

\hypertarget{ux4e3aux4ec0ux4e48ux8981ux727aux7272ux4e00ux4e2aux5b58ux50a8ux5355ux5143}{%
\subsection{为什么要牺牲一个存储单元}\label{ux4e3aux4ec0ux4e48ux8981ux727aux7272ux4e00ux4e2aux5b58ux50a8ux5355ux5143}}

当队列只剩最后一个空闲空间时，就认为\hlred{队列已经满了}。为什么不把最后这个位置也用上？因为初始化时让
\texttt{front} 和 \texttt{rear}
指向同一个位置，并且用“两个指针是否指向同一位置”来判断队空。如果把最后一个空位也填上，\texttt{rear}
后移后会和 \texttt{front} 重合，按刚才的逻辑这就成了“队空”，显然不对。

所以必须\uln{牺牲一个存储单元}：入队前检查队尾指针的\textbf{再下一个位置}是不是队头。

\begin{defbox}

\textbf{队满条件}（rear
指向队尾元素后一个位置，牺牲一个存储单元）：\texttt{(Q.rear\ +\ 1)\ \%\ MaxSize\ ==\ Q.front}。

\end{defbox}

\slidefig{assets/3.2.2-p07.jpg}{0.6}{课件：循环队列；队满条件与代价}

\begin{Shaded}
\begin{Highlighting}[]
\CommentTok{// 入队}
\DataTypeTok{bool}\NormalTok{ EnQueue}\OperatorTok{(}\NormalTok{SqQueue }\OperatorTok{\&}\NormalTok{Q}\OperatorTok{,}\NormalTok{ ElemType x}\OperatorTok{)} \OperatorTok{\{}
    \ControlFlowTok{if} \OperatorTok{((}\NormalTok{Q}\OperatorTok{.}\NormalTok{rear }\OperatorTok{+} \DecValTok{1}\OperatorTok{)} \OperatorTok{\%}\NormalTok{ MaxSize }\OperatorTok{==}\NormalTok{ Q}\OperatorTok{.}\NormalTok{front}\OperatorTok{)}   \CommentTok{// 判断队满}
        \ControlFlowTok{return} \KeywordTok{false}\OperatorTok{;}                        \CommentTok{// 队满则报错}
\NormalTok{    Q}\OperatorTok{.}\NormalTok{data}\OperatorTok{[}\NormalTok{Q}\OperatorTok{.}\NormalTok{rear}\OperatorTok{]} \OperatorTok{=}\NormalTok{ x}\OperatorTok{;}                      \CommentTok{// 新元素插入队尾}
\NormalTok{    Q}\OperatorTok{.}\NormalTok{rear }\OperatorTok{=} \OperatorTok{(}\NormalTok{Q}\OperatorTok{.}\NormalTok{rear }\OperatorTok{+} \DecValTok{1}\OperatorTok{)} \OperatorTok{\%}\NormalTok{ MaxSize}\OperatorTok{;}         \CommentTok{// 队尾指针加 1 取模}
    \ControlFlowTok{return} \KeywordTok{true}\OperatorTok{;}
\OperatorTok{\}}
\end{Highlighting}
\end{Shaded}

\slidefig{assets/3.2.2-p08.jpg}{0.6}{课件：循环队列的入队操作（判断队空、判断队满）}

\hypertarget{ux56dbux51faux961fux4e0eux8bfbux961fux5934}{%
\section{四、出队与读队头}\label{ux56dbux51faux961fux4e0eux8bfbux961fux5934}}

\tocback

\Needspace{4\baselineskip}

\cuenote{出队时 \texttt{front} 怎么移动？\texttt{GetHead} 比 \texttt{DeQueue}
少了哪一句？}只能让队头元素出队。先判空，空则返回 \texttt{false}；否则把
\texttt{front} 指向的元素赋给 \texttt{x}，再让 \texttt{front}
后移一位，\uln{同样要对 \texttt{MaxSize} 取模}，\texttt{front}
才能转着圈移动。一直出队，直到 \texttt{front} 和 \texttt{rear}
再次指向同一位置，队列被取空。

\begin{Shaded}
\begin{Highlighting}[]
\CommentTok{// 出队（删除一个队头元素，并用 x 返回）}
\DataTypeTok{bool}\NormalTok{ DeQueue}\OperatorTok{(}\NormalTok{SqQueue }\OperatorTok{\&}\NormalTok{Q}\OperatorTok{,}\NormalTok{ ElemType }\OperatorTok{\&}\NormalTok{x}\OperatorTok{)} \OperatorTok{\{}
    \ControlFlowTok{if} \OperatorTok{(}\NormalTok{Q}\OperatorTok{.}\NormalTok{rear }\OperatorTok{==}\NormalTok{ Q}\OperatorTok{.}\NormalTok{front}\OperatorTok{)}                   \CommentTok{// 判断队空}
        \ControlFlowTok{return} \KeywordTok{false}\OperatorTok{;}                        \CommentTok{// 队空则报错}
\NormalTok{    x }\OperatorTok{=}\NormalTok{ Q}\OperatorTok{.}\NormalTok{data}\OperatorTok{[}\NormalTok{Q}\OperatorTok{.}\NormalTok{front}\OperatorTok{];}
\NormalTok{    Q}\OperatorTok{.}\NormalTok{front }\OperatorTok{=} \OperatorTok{(}\NormalTok{Q}\OperatorTok{.}\NormalTok{front }\OperatorTok{+} \DecValTok{1}\OperatorTok{)} \OperatorTok{\%}\NormalTok{ MaxSize}\OperatorTok{;}       \CommentTok{// 队头指针后移}
    \ControlFlowTok{return} \KeywordTok{true}\OperatorTok{;}
\OperatorTok{\}}

\CommentTok{// 获得队头元素的值，用 x 返回}
\DataTypeTok{bool}\NormalTok{ GetHead}\OperatorTok{(}\NormalTok{SqQueue Q}\OperatorTok{,}\NormalTok{ ElemType }\OperatorTok{\&}\NormalTok{x}\OperatorTok{)} \OperatorTok{\{}
    \ControlFlowTok{if} \OperatorTok{(}\NormalTok{Q}\OperatorTok{.}\NormalTok{rear }\OperatorTok{==}\NormalTok{ Q}\OperatorTok{.}\NormalTok{front}\OperatorTok{)}
        \ControlFlowTok{return} \KeywordTok{false}\OperatorTok{;}                        \CommentTok{// 队空则报错}
\NormalTok{    x }\OperatorTok{=}\NormalTok{ Q}\OperatorTok{.}\NormalTok{data}\OperatorTok{[}\NormalTok{Q}\OperatorTok{.}\NormalTok{front}\OperatorTok{];}
    \ControlFlowTok{return} \KeywordTok{true}\OperatorTok{;}
\OperatorTok{\}}
\end{Highlighting}
\end{Shaded}

队列里通常只需要读队头元素。\texttt{GetHead} 就是把出队里“\texttt{front}
后移”那一句删掉。

\slidefig{assets/3.2.2-p09.jpg}{0.6}{课件：循环队列的出队操作与读队头}

\hypertarget{ux4e94ux5224ux7a7aux5224ux6ee1ux7684ux4e09ux79cdux65b9ux6848}{%
\section{五、判空判满的三种方案}\label{ux4e94ux5224ux7a7aux5224ux6ee1ux7684ux4e09ux79cdux65b9ux6848}}

\tocback

\cueblock{\cuenote}{\relax{三种方案各增加了什么、各自的队空、队满条件是什么？队列元素个数怎么算？}}

\hypertarget{ux65b9ux6848ux4e00ux727aux7272ux4e00ux4e2aux5b58ux50a8ux5355ux5143}{%
\subsection{方案一：牺牲一个存储单元}\label{ux65b9ux6848ux4e00ux727aux7272ux4e00ux4e2aux5b58ux50a8ux5355ux5143}}

初始化 \texttt{rear\ =\ front\ =\ 0}。

\begin{itemize}
\tightlist
\item
  队空：\texttt{Q.rear\ ==\ Q.front}；
\item
  队满：\texttt{(Q.rear\ +\ 1)\ \%\ MaxSize\ ==\ Q.front}；
\item
  \textbf{队列元素个数}：\hlred{\texttt{(rear\ +\ MaxSize\ -\ front)\ \%\ MaxSize}}。
\end{itemize}

课件例子：\texttt{rear\ =\ 2}，\texttt{front\ =\ 3}，\texttt{MaxSize\ =\ 10}，元素个数
\(= (2+10-3)\%10 = 9\)。老师说：学过数论的同学可以用模运算的性质证明这个式子；不会证明的也要记住，\hlred{这个挺常考}，可以多找例子验证。

\slidefig{assets/3.2.2-p10.jpg}{0.6}{课件：方案一（“出题老师：不准浪费这可怜孤独的存储空间！！！”）}

\begin{addbox}

\textbf{元素个数公式为什么对}（我的说明）：不绕圈时
\texttt{rear\ ≥\ front}，个数就是 \texttt{rear\ -\ front}，加
\texttt{MaxSize} 再取模不改变结果；绕了圈时
\texttt{rear\ \textless{}\ front}，\texttt{rear\ -\ front}
是负数，真实个数是
\texttt{rear\ -\ front\ +\ MaxSize}（把绕过去的一圈补回来），它在
0～\texttt{MaxSize}−1 之间，取模也不改变它。加 \texttt{MaxSize}
是为了让结果不为负，取模是为了不绕圈时把多加的 \texttt{MaxSize} 去掉。

\end{addbox}

方案一以牺牲一片存储空间为代价，自己写项目已经完全够用。但有的题目会要求\uln{不能浪费这片存储空间}，而把它也存满后，队满和队空的状态一样（都是
\texttt{front\ ==\ rear}），代码区分不了。下面两种方案解决这个问题。

\hypertarget{ux65b9ux6848ux4e8cux589eux52a0-size-ux53d8ux91cf}{%
\subsection{方案二：增加 size
变量}\label{ux65b9ux6848ux4e8cux589eux52a0-size-ux53d8ux91cf}}

在结构体里定义 \texttt{int\ size;} 记录队列当前长度。

\begin{itemize}
\tightlist
\item
  初始化：\texttt{rear\ =\ front\ =\ 0;\ size\ =\ 0;}
\item
  插入成功 \texttt{size++}，删除成功 \texttt{size-\/-}；
\item
  队空：\texttt{size\ ==\ 0}；队满：\texttt{size\ ==\ MaxSize}；元素个数就是
  \texttt{size}。
\end{itemize}

虽然队空和队满时 \texttt{front}、\texttt{rear} 都指向同一位置，但可以用
\texttt{size} 区分。

\slidefig{assets/3.2.2-p11.jpg}{0.6}{课件：方案二（增加 size）}

\begin{fixbox}

正课字幕 P24 这里是“每一个新元素入队成功，我们都需要让 size
的值减减，每一个对头元素出对成功”，漏了半句。课件写的是\uln{插入成功 \texttt{size++}；删除成功 \texttt{size-\/-}}，以课件为准。

\end{fixbox}

\hypertarget{ux65b9ux6848ux4e09ux589eux52a0-tag-ux53d8ux91cf}{%
\subsection{方案三：增加 tag
变量}\label{ux65b9ux6848ux4e09ux589eux52a0-tag-ux53d8ux91cf}}

在结构体里定义 \texttt{int\ tag;} 记录最近进行的是删除还是插入：

\begin{itemize}
\tightlist
\item
  每次删除操作成功时，都令 \texttt{tag\ =\ 0}；每次插入操作成功时，都令
  \texttt{tag\ =\ 1}；
\item
  初始化：\texttt{rear\ =\ front\ =\ 0;\ tag\ =\ 0;}
\item
  \uln{只有删除操作才可能导致队空，只有插入操作才可能导致队满}；
\item
  队空：\texttt{front\ ==\ rear\ \&\&\ tag\ ==\ 0}；
\item
  队满：\texttt{front\ ==\ rear\ \&\&\ tag\ ==\ 1}。
\end{itemize}

例：只剩最后一个空位，不管此时 \texttt{tag} 是多少，插入成功后
\texttt{tag} 都置 1，\texttt{rear} 后移和 \texttt{front}
重合，这次重合是插入导致的，所以是队满。只剩最后一个元素，删除后
\texttt{tag} 置 0，\texttt{front} 后移和 \texttt{rear}
重合，是删除导致的，所以是队空。

\slidefig{assets/3.2.2-p12.jpg}{0.6}{课件：方案三（增加 tag）}

\begin{addbox}

课件这一页在左上角的元素个数公式 \texttt{(rear+MaxSize-front)\%MaxSize}
旁画了一个红叉，老师没有口头解释。我的理解（推断）：方案三队满时
\texttt{front\ ==\ rear}，公式算出来是 0 而不是
\texttt{MaxSize}，所以\uln{方案三算长度要先看 \texttt{tag}}：队满时长度是
\texttt{MaxSize}，其他情况才用公式。

\end{addbox}

\begin{bookbox}

王道书第 77--78 页“循环队列”：

\begin{itemize}
\tightlist
\item
  初始状态 \texttt{Q.front=Q.rear=0}；队首指针进
  1：\texttt{Q.front=(Q.front+1)\%MaxSize}；队尾指针进
  1：\texttt{Q.rear=(Q.rear+1)\%MaxSize}；
\item
  队列长度：\texttt{(Q.rear+MaxSize-Q.front)\%MaxSize}；
\item
  出入队时：指针都按\uln{顺时针方向}移动（见图 3.7）。
\end{itemize}

队空条件显然是
\texttt{Q.front==Q.rear}。但当入队速度远快于出队时，\texttt{rear} 会追上
\texttt{front}，此时同样满足
\texttt{Q.front==Q.rear}，却表示队满。因此仅凭 \texttt{front==rear}
无法区分队空与队满。常用以下三种方法：

\begin{enumerate}
\def\labelenumi{\arabic{enumi}.}
\tightlist
\item
  \textbf{牺牲一个存储单元}（较为普遍的做法）：约定以“队尾指针的下一位置为队首”作为队满标志。队满
  \texttt{(Q.rear+1)\%MaxSize==Q.front}；队空
  \texttt{Q.front==Q.rear}；队列中元素个数
  \texttt{(Q.rear-Q.front+MaxSize)\%MaxSize}。
\item
  \textbf{增设 size 成员}：入队成功 \texttt{size++}，出队成功
  \texttt{size-\/-}。队空 \texttt{Q.size==0}，队满
  \texttt{Q.size==MaxSize}。两种情形均有 \texttt{Q.front==Q.rear}，由
  \texttt{size} 区分。
\item
  \textbf{增设 tag 标志位}：出队后置 \texttt{tag=0}，若此时
  \texttt{Q.front==Q.rear}，则为队空；入队后置 \texttt{tag=1}，若此时
  \texttt{Q.front==Q.rear}，则为队满。
\end{enumerate}

书上第 78--79
页“循环队列的操作（基于牺牲一个单元法）”的初始化、判队空、入队、出队代码和课件相同（书上
\texttt{MaxSize} 是
50）。考点追踪：特定条件下循环队列队头/队尾指针的初值（2011）；循环队列判空/满的条件（2014）。

\end{bookbox}

\slidefig{assets/3.2-book-fig3.7.png}{0.6}{王道书图 3.7：循环队列出入队的示意图；(d1) 无法判断队满还是队空，(d2)
牺牲一个存储单元}

\begin{addbox}

\textbf{名词统一}：书上队列长度写作
\texttt{(Q.rear-Q.front+MaxSize)\%MaxSize}，课件写作
\texttt{(rear+MaxSize-front)\%MaxSize}，只是加法交换了顺序，同一个式子。

\end{addbox}

\hypertarget{ux516dux5176ux4ed6ux51faux9898ux65b9ux6cd5rear-ux6307ux5411ux961fux5c3eux5143ux7d20}{%
\section{六、其他出题方法：rear
指向队尾元素}\label{ux516dux5176ux4ed6ux51faux9898ux65b9ux6cd5rear-ux6307ux5411ux961fux5c3eux5143ux7d20}}

\tocback

\Needspace{4\baselineskip}

\cuenote{\texttt{rear} 指向队尾元素时，入队的两句顺序怎么变？初始化时
\texttt{rear} 应该指向哪里？判空、判满各是什么？}前面所有方法都基于“队尾指针指向队尾元素的\textbf{下一个}位置”。考试中也可能遇到\hlred{队尾指针\textbf{指向队尾元素}}的情况，代码会有不同。考试时一定要注意审题。

\textbf{入队}：要先让 \texttt{rear} 后移一位，再往这个位置写入
\texttt{x}：

\begin{Shaded}
\begin{Highlighting}[]
\NormalTok{Q}\OperatorTok{.}\NormalTok{rear }\OperatorTok{=} \OperatorTok{(}\NormalTok{Q}\OperatorTok{.}\NormalTok{rear }\OperatorTok{+} \DecValTok{1}\OperatorTok{)} \OperatorTok{\%}\NormalTok{ MaxSize}\OperatorTok{;}    \CommentTok{// rear 指向队尾元素：先后移}
\NormalTok{Q}\OperatorTok{.}\NormalTok{data}\OperatorTok{[}\NormalTok{Q}\OperatorTok{.}\NormalTok{rear}\OperatorTok{]} \OperatorTok{=}\NormalTok{ x}\OperatorTok{;}                 \CommentTok{// 再写入}
\end{Highlighting}
\end{Shaded}

出队不变：\texttt{x\ =\ Q.data{[}Q.front{]};\ Q.front\ =\ (Q.front\ +\ 1)\ \%\ MaxSize;}

\textbf{初始化}：因为每次入队都是先让 \texttt{rear}
后移再放元素，比较合理的方式是\uln{\texttt{front} 指向 0，\texttt{rear} 指向 \texttt{n-1}}（\texttt{MaxSize\ -\ 1}）。这样插入第一个元素时
\texttt{rear} 先后移到 0，再把 \texttt{x} 放进 0
号位置。老师说这种初始化方式可能不太符合直觉。

\slidefig{assets/3.2.2-p13.jpg}{0.6}{课件：其他出题方法（初始时队尾指针指向哪更合理？）}

\textbf{判空}：队尾指针的下一个位置是队头：\texttt{(Q.rear\ +\ 1)\ \%\ MaxSize\ ==\ Q.front}。

\textbf{判满}：如果所有位置都填满，\texttt{rear} 的下一个位置也是
\texttt{front}，和队空一样，区分不了。处理方法和之前一样：

\begin{itemize}
\tightlist
\item
  方案一：规定队头指针前面的这个存储单元不能存数据（牺牲一个存储单元），\texttt{rear}
  最多只能走到它前面。于是\uln{头指针在尾指针后一个位置是队空，在尾指针后两个位置是队满}；
\item
  方案二：增加辅助变量 \texttt{size} 或 \texttt{tag}，这里不再展开。
\end{itemize}

\slidefig{assets/3.2.2-p15.jpg}{0.6}{课件：rear
指向队尾元素时的判空与判满（“判满：(Q.rear+1)\%MaxSize==Q.front？”打叉）}

\begin{bookbox}

王道书第 76 页：队首指针 \texttt{front} 指向队首元素，队尾指针
\texttt{rear}
指向队尾元素的下一个位置。\uln{不同教材对 \texttt{front} 和 \texttt{rear} 指针含义的约定可能不同}（例如，\texttt{rear}
可能指向队尾元素或其下一位置），这会导致入队和出队操作的具体实现不同。

\end{bookbox}

\begin{addbox}

\textbf{老师留的思考题}：“如果让队尾指针指向队尾元素，同时采用增加
\texttt{size}
变量记录队列长度的策略，这些基本操作怎么实现？”课件知识回顾里把它扩展成
①a、①b、①c、②a、②b、②c 六种组合（① \texttt{rear}
指向队尾元素后一个位置；② \texttt{rear} 指向队尾元素；a
牺牲一个存储单元；b 增加 \texttt{size}；c 增加
\texttt{tag}）。建议先自己想，再对照我的答案。记 \(M\) =
\texttt{MaxSize}：

\begin{itemize}
\tightlist
\item
  \textbf{初始化}：① 类 \texttt{front\ =\ rear\ =\ 0}；② 类
  \texttt{front\ =\ 0;\ rear\ =\ M-1}。b 方案再加 \texttt{size\ =\ 0}，c
  方案再加 \texttt{tag\ =\ 0}。
\end{itemize}

{\rowcolors{2}{wdblue!5}{white}\arrayrulecolor{wdblue!35}
\begin{xltabular}{\linewidth}{>{\hsize=0.338\hsize\raggedright\arraybackslash}X>{\hsize=1.183\hsize\raggedright\arraybackslash}X>{\hsize=1.183\hsize\raggedright\arraybackslash}X>{\hsize=1.296\hsize\raggedright\arraybackslash}X}
\rowcolor{wdblue}\textcolor{white}{\bfseries 方案} & \textcolor{white}{\bfseries 判空} & \textcolor{white}{\bfseries 判满} & \textcolor{white}{\bfseries 队列长度} \\
\endfirsthead
①a & \texttt{rear\ ==\ front} & \texttt{(rear+1)\%M\ ==\ front} & \texttt{(rear-front+M)\%M} \\ \hline
①b & \texttt{size\ ==\ 0} & \texttt{size\ ==\ M} & \texttt{size} \\ \hline
①c & \texttt{rear\ ==\ front\ \&\&\ tag\ ==\ 0} & \texttt{rear\ ==\ front\ \&\&\ tag\ ==\ 1} & 队满时 \(M\)，否则同 ①a \\ \hline
②a & \texttt{(rear+1)\%M\ ==\ front} & \texttt{(rear+2)\%M\ ==\ front} & \texttt{(rear-front+1+M)\%M} \\ \hline
②b & \texttt{size\ ==\ 0} & \texttt{size\ ==\ M} & \texttt{size} \\ \hline
②c & \texttt{(rear+1)\%M\ ==\ front\ \&\&\ tag\ ==\ 0} & \texttt{(rear+1)\%M\ ==\ front\ \&\&\ tag\ ==\ 1} & 队满时 \(M\)，否则同 ②a \\ \hline
\end{xltabular}\par\vspace{4pt}}

\begin{itemize}
\tightlist
\item
  入队：① 先放再移（\texttt{data{[}rear{]}=x;\ rear=(rear+1)\%M;}）；②
  先移再放（\texttt{rear=(rear+1)\%M;\ data{[}rear{]}=x;}）。出队六种都一样：\texttt{x=data{[}front{]};\ front=(front+1)\%M;}
\item
  b 方案入队成功 \texttt{size++}、出队成功 \texttt{size-\/-}；c
  方案入队成功 \texttt{tag=1}、出队成功 \texttt{tag=0}；
\item
  a 方案最多存 \(M-1\) 个元素，b、c 方案能存满 \(M\) 个。
\end{itemize}

这张表我用程序模拟验证过：\(M=6\)，六种方案各随机做 3000 轮、每轮 40
次入队/出队，每一步的判空、判满、长度都和真实队列一致。

\end{addbox}

\hypertarget{ux4e03ux77e5ux8bc6ux56deux987e}{%
\section{七、知识回顾}\label{ux4e03ux77e5ux8bc6ux56deux987e}}

\tocback

\slidefig{assets/3.2.2-p16.jpg}{0.65}{课件：知识回顾与重要考点（思考：分别采用 ①a、①b、①c、②a、②b、②c
策略时，这些操作怎么实现？）}

老师的总结：这一节很容易在选择题中考察，要注意两个方面：

\begin{enumerate}
\def\labelenumi{\arabic{enumi}.}
\tightlist
\item
  队尾指针到底指向\uln{队尾元素的后一个位置}，还是\uln{指向队尾元素}；
\item
  题目给的判空、判满方法是什么：牺牲一个存储单元，还是增加辅助变量（\texttt{size}、\texttt{tag}）。
\end{enumerate}

把这几种情况想明白，这一节的考题基本没有什么难得住的。

\begin{recallbox}

正课 3.1.2 讲顺序栈时也有同样的“先审题”：\texttt{top}
指向栈顶元素还是栈顶元素的下一个位置，两种设计的进栈、出栈顺序相反。这里的
\texttt{rear}
是同一个道理：\uln{指向空位就先放再移，指向元素就先移再放}。

\end{recallbox}

\begin{summarybox}

\textbf{顺序队列} = 静态数组 \texttt{data{[}MaxSize{]}} +
\texttt{front}（指向队头元素）+
\texttt{rear}（默认指向队尾元素的后一个位置）。

\textbf{循环队列}：指针后移都写成
\texttt{(指针+1)\%MaxSize}，用模运算把存储空间在逻辑上变成“环状”，解决
\texttt{rear==MaxSize}
时数组前面还有空位的问题（书上叫\textbf{假溢出}）。

\textbf{rear 指向队尾元素后一个位置时}：

\begin{itemize}
\tightlist
\item
  初始化 \texttt{front=rear=0}；入队
  \texttt{data{[}rear{]}=x;\ rear=(rear+1)\%MaxSize}；出队
  \texttt{x=data{[}front{]};\ front=(front+1)\%MaxSize}；
\item
  方案一（牺牲一个单元）：队空 \texttt{rear==front}，队满
  \hlred{\texttt{(rear+1)\%MaxSize==front}}；
\item
  方案二（\texttt{size}）：队空 \texttt{size==0}，队满
  \texttt{size==MaxSize}；
\item
  方案三（\texttt{tag}，删除置 0、插入置 1）：队空
  \texttt{front==rear\ \&\&\ tag==0}，队满
  \texttt{front==rear\ \&\&\ tag==1}；
\item
  元素个数 \hlred{\texttt{(rear+MaxSize-front)\%MaxSize}}（常考）。
\end{itemize}

\textbf{rear 指向队尾元素时}：初始化
\texttt{front=0,\ rear=MaxSize-1}；入队先移再放；判空
\texttt{(rear+1)\%MaxSize==front}；牺牲一个单元时判满
\texttt{(rear+2)\%MaxSize==front}。

\textbf{自测}：① 为什么 \texttt{rear==MaxSize} 不能作为队满条件？②
牺牲一个单元时队满条件是什么，为什么要牺牲？③
\texttt{MaxSize=10}，\texttt{front=7}，\texttt{rear=3}（rear
指向队尾元素后一个位置），队中有几个元素？④ \texttt{tag} 方案中，为什么
\texttt{front==rear\ \&\&\ tag==1} 一定是队满？⑤ \texttt{rear}
指向队尾元素时，初始化 \texttt{rear} 设成多少，入队两句怎么写？

\end{summarybox}
\end{document}
